% GENERATED FILE. Edit canonical lessons, then run npm run generate:books.
% canonical-source-hash: fnv1a64:177922d825c7ae04
% canonical-lessons: FR-C429-ampoule, FR-C429-prise, FR-C429-minibar, FR-C429-moquette, FR-C429-tache-stain

\chapter{Petits soucis dans la chambre}
\label{ch:fr-petits-soucis-dans-la-chambre}
\hlchaptermodality{429}

\begin{chapteropening}
\textbf{By the end of this chapter:} \emph{I can describe small problems in a room: the light bulb, the socket, the minibar, the carpet and a stain.}

Complete the last lesson of chapter 429: I can describe small problems in a room: the light bulb, the socket, the minibar, the carpet and a stain.
\end{chapteropening}

\section[l'ampoule]{l'ampoule — the light bulb; the blister}

\label{lesson:FR-C429-ampoule}



\begin{quote}
Before the new one: say the French for to drip, then the French for to unblock.
\end{quote}

\subsection*{You'll want to know: l'ampoule}
\textbf{l'ampoule} — \textquotedblleft{}the light bulb; the blister\textquotedblright{}.

\begin{quote}
L'ampoule de la lampe ne marche plus.
\end{quote}

\begin{quote}
\emph{The bulb in the lamp has gone.}
\end{quote}

A feminine noun. After a long walk, \emph{une ampoule} is a blister.

\begin{cousinweb}
\emph{ampoule} is inherited from Latin \emph{ampulla} 'small flask', built on \emph{amphora} (EtymDB): a bulb looks like a little glass flask.
\end{cousinweb}

\begin{grammarlens}[title={What you've built}]
One more word for telling, explaining and describing at B1.
\end{grammarlens}

\subsection*{Your turn}
\begin{itemize}
\raggedright
  \item \emph{Say it:} \emph{l'ampoule}
  \item \emph{Say it:} \emph{l'ampoule}, once more
  \item \emph{Say it:} a sentence of your own with \emph{l'ampoule}
  \item \emph{Recall:} say the French for to drip, then the French for to unblock, then say \emph{l'ampoule} again
\end{itemize}

\begin{tcolorbox}[breakable,colback=teal!4,colframe=teal!35!black,title={Before you move on}]
What does l'ampoule mean? (\textquotedblleft{}The light bulb; the blister\textquotedblright{}.) Say it once more.
\end{tcolorbox}
\section[la prise]{la prise — the socket (for electricity), the plug; the hold}

\label{lesson:FR-C429-prise}



\begin{quote}
Before the new one: say the French for to unblock, then the French for the light bulb.
\end{quote}

\subsection*{You'll want to know: la prise}
\textbf{la prise} — \textquotedblleft{}the socket (for electricity), the plug; the hold\textquotedblright{}.

\begin{quote}
Il n'y a pas de prise près du lit.
\end{quote}

\begin{quote}
\emph{There is no socket near the bed.}
\end{quote}

\emph{Une prise de courant} — a socket; \emph{une prise} also a grip or hold.

\begin{cousinweb}
\emph{prise} is the feminine past participle of \emph{prendre} 'to take', used as a noun (EtymDB); \emph{prendre} is inherited from Latin \emph{prehendere} 'to grasp'.
\end{cousinweb}

\begin{grammarlens}[title={What you've built}]
One more word for telling, explaining and describing at B1.
\end{grammarlens}

\subsection*{Your turn}
\begin{itemize}
\raggedright
  \item \emph{Say it:} \emph{la prise}
  \item \emph{Say it:} \emph{la prise}, once more
  \item \emph{Say it:} a sentence of your own with \emph{la prise}
  \item \emph{Recall:} say the French for to unblock, then the French for the light bulb, then say \emph{la prise} again
\end{itemize}

\begin{tcolorbox}[breakable,colback=teal!4,colframe=teal!35!black,title={Before you move on}]
What does la prise mean? (\textquotedblleft{}The socket (for electricity), the plug; the hold\textquotedblright{}.) Say it once more.
\end{tcolorbox}
\section[le minibar]{le minibar — the minibar}

\label{lesson:FR-C429-minibar}



\begin{quote}
Before the new one: say the French for the light bulb, then the French for the socket.
\end{quote}

\subsection*{You'll want to know: le minibar}
\textbf{le minibar} — \textquotedblleft{}the minibar\textquotedblright{}.

\begin{quote}
Les boissons du minibar sont très chères.
\end{quote}

\begin{quote}
\emph{The drinks in the minibar are very dear.}
\end{quote}

Plural \emph{les minibars}.

\begin{cousinweb}
\emph{minibar} is built from \emph{mini-} 'small' and \emph{bar}, which was borrowed from English \emph{bar}. EtymDB does not list it.
\end{cousinweb}

\begin{grammarlens}[title={What you've built}]
One more word for telling, explaining and describing at B1.
\end{grammarlens}

\subsection*{Your turn}
\begin{itemize}
\raggedright
  \item \emph{Say it:} \emph{le minibar}
  \item \emph{Say it:} \emph{le minibar}, once more
  \item \emph{Say it:} a sentence of your own with \emph{le minibar}
  \item \emph{Recall:} say the French for the light bulb, then the French for the socket, then say \emph{le minibar} again
\end{itemize}

\begin{tcolorbox}[breakable,colback=teal!4,colframe=teal!35!black,title={Before you move on}]
What does le minibar mean? (\textquotedblleft{}The minibar\textquotedblright{}.) Say it once more.
\end{tcolorbox}
\section[la moquette]{la moquette — the fitted carpet}

\label{lesson:FR-C429-moquette}



\begin{quote}
Before the new one: say the French for the socket, then the French for the minibar.
\end{quote}

\subsection*{You'll want to know: la moquette}
\textbf{la moquette} — \textquotedblleft{}the fitted carpet\textquotedblright{}.

\begin{quote}
La moquette de la chambre est un peu sale.
\end{quote}

\begin{quote}
\emph{The carpet in the room is a bit dirty.}
\end{quote}

\emph{Un tapis} is a rug; \emph{la moquette} covers the whole floor. In Belgium also \emph{le tapis plain}.

\begin{cousinweb}
The origin of \emph{moquette} is not settled; it was first a kind of thick cloth. EtymDB lists it without a source.
\end{cousinweb}

\begin{grammarlens}[title={What you've built}]
One more word for telling, explaining and describing at B1.
\end{grammarlens}

\subsection*{Your turn}
\begin{itemize}
\raggedright
  \item \emph{Say it:} \emph{la moquette}
  \item \emph{Say it:} \emph{la moquette}, once more
  \item \emph{Say it:} a sentence of your own with \emph{la moquette}
  \item \emph{Recall:} say the French for the socket, then the French for the minibar, then say \emph{la moquette} again
\end{itemize}

\begin{tcolorbox}[breakable,colback=teal!4,colframe=teal!35!black,title={Before you move on}]
What does la moquette mean? (\textquotedblleft{}The fitted carpet\textquotedblright{}.) Say it once more.
\end{tcolorbox}
\section[la tache]{la tache — the stain, the spot}

\label{lesson:FR-C429-tache-stain}



\begin{quote}
Before the new one: say the French for the minibar, then the French for the fitted carpet.
\end{quote}

\subsection*{You'll want to know: la tache}
\textbf{la tache} — \textquotedblleft{}the stain, the spot\textquotedblright{}.

\begin{quote}
Il y a une tache de café sur le drap.
\end{quote}

\begin{quote}
\emph{There is a coffee stain on the sheet.}
\end{quote}

No accent: \emph{une tache} is a stain; \emph{une tâche} is a task.

\begin{cousinweb}
\emph{tache} comes from Old French \emph{tache} (EtymDB), which is usually traced to a Germanic word meaning 'sign, mark'.
\end{cousinweb}

\begin{grammarlens}[title={What you've built}]
One more word for telling, explaining and describing at B1.
\end{grammarlens}

\subsection*{Your turn}
\begin{itemize}
\raggedright
  \item \emph{Say it:} \emph{la tache}
  \item \emph{Say it:} \emph{la tache}, once more
  \item \emph{Say it:} a sentence of your own with \emph{la tache}
  \item \emph{Recall:} say the French for the minibar, then the French for the fitted carpet, then say \emph{la tache} again
\end{itemize}

\begin{tcolorbox}[breakable,colback=teal!4,colframe=teal!35!black,title={Before you move on}]
What does la tache mean? (\textquotedblleft{}The stain, the spot\textquotedblright{}.) Say it once more.
\end{tcolorbox}
